\scanpage{158}
\uline{Involutions.} Supposons $\alpha\in B(M,M)$ forme \deleted{\unsure{quelconque}} \deleted{\unsure{et}} non dégénérée, et sym. ou \deleted{anti} alternée, alors \deleted{Alors par} \unsure{l'iso} $\widetilde\alpha:M\simeq M^{*}$. Soit de même $\beta\in B(N,N)$ d'où $\widetilde\beta:N\longrightarrow N^{*}$ $(M,N\text{ de deg }i)$.
Soit $u:M\mathbin{-}N$, on définit alors $u':N\mathbin{-}M$ par comm. du diagramme ci-contre.
\diagram{\begin{tikzcd}
M \arrow[r,leftrightarrow,"u^{\prime\deleted{\unsure{*}}}"] \arrow[d,"\widetilde\alpha"'] & N \arrow[d,"\widetilde\beta"]\\
M^{*} & N^{*} \arrow[l,"\unsure{{}^tu^{*}|u^{*}}"']
\end{tikzcd}}
\margin{transposé}
On trouve les relations habituelles
\[
\begin{aligned}
(u\pm v)'&=u'\pm v'\\
(uv)'&=v'u'\\
1'&=1\\
(u')'&=u
\end{aligned}
\]
\margin{si $\alpha$ et $\beta$ sont tous deux \deleted{\unsure{sym.}} \inserted{sym.} ou tous deux alternés.}
En particulier sur $\operatorname{End}(M)$ on a une structure d'algèbre involutive sur $\mathbf Q$. \deleted{\unsure{Alors}}
Si $\alpha$ est remplacé par $\alpha_1$, $\beta$ par $\beta_1$, on aura
\[
\widetilde\alpha_1=\widetilde\alpha\circ A,\qquad\widetilde\beta_1=\widetilde\beta\circ B\qquad(A\in\operatorname{End}(M),\ B\in\operatorname{End}(N))
\]
et la nouvelle ``adjonction'' sera
\[
u''=\deleted{\unsure{\widetilde\beta_1}}\ \widetilde\alpha_1^{-1}\circ u^{*}\circ\widetilde\beta_1
=A^{-1}\widetilde\alpha^{-1}u^{*}\widetilde\beta B\qquad\text{i.e}
\]
\[
u''=A^{-1}u'B
\]
En particulier, si $M=N$, on a dans $A$ une involution définie mod \deleted{\unsure{autom. intérieurs}}\inserted{\unsure{transformation} de $u\mapsto u'$ en} $u\mapsto A^{-1}u'A$, où $A$ est un élément autoadjoint (?)
\mspage{157}
